% Solutions to the exercises of Chapter 12.
\begin{solutions}
  \item Let the prior be $(\hat x, P)$ and the measurements $z_i = x + n_i$ with variances $R_i$. \emph{Sequentially}: the first update gives
        $K_1 = P/(P + R_1)$, $\hat x_1 = \hat x + K_1(z_1 - \hat x)$ and $P_1 = PR_1/(P + R_1)$, i.e.\ $P_1\inv = P\inv + R_1\inv$ and $P_1\inv\hat x_1 = P\inv\hat x + R_1\inv z_1$; the second
        adds $R_2\inv$ and $R_2\inv z_2$ in the same way:
        \[
          P_2\inv = P\inv + R_1\inv + R_2\inv, \qquad \hat x_2 = P_2\big(P\inv\hat x + R_1\inv z_1 + R_2\inv z_2\big).
        \]
        \emph{Stacked}: $H = (1, 1)\T$ and $R = \diag(R_1, R_2)$. Then $S = \begin{bmatrix}P + R_1 & P\\ P & P + R_2\end{bmatrix}$ with $\det S = P(R_1 + R_2) + R_1R_2$, and
        $K = P(1, 1)S\inv = \frac{P}{\det S}(R_2, R_1)$. Hence $P^+ = P - KHP = P - \frac{P^2(R_1 + R_2)}{\det S} = \frac{PR_1R_2}{\det S}$, whose inverse is $P\inv + R_1\inv + R_2\inv$, and
        \[
          \hat x^+ = \hat x + \frac{P}{\det S}\big(R_2(z_1 - \hat x) + R_1(z_2 - \hat x)\big),
        \]
        which, multiplied out, equals $\hat x_2$. Both routes give the same posterior, as the sequential-processing proposition states.

  \item By \eqref{eq:kf-update}, $P^- - P^+ = KSK\T$, which is PSD because $S \succ 0$ (it is a congruence of $S$); so $P^+ \preceq P^-$. For an arbitrary gain, write
        $K = K_\star + D$ with $K_\star$ optimal. Expanding the Joseph form and using $K_\star S = P^-H\T$,
        \[
          P^+(K) = (I - KH)P^-(I - KH)\T + KRK\T = P^+(K_\star) + DSD\T .
        \]
        The optimal gain therefore minimizes the posterior covariance, but a poor gain can make it \emph{larger} than the prior: for $P^- = H = R = 1$ and
        $K = 2$, $P^+ = (1 - 2)^2 + 4 = 5 > 1$. That is correct, not a flaw: a badly weighted update really does make the estimate worse, and the Joseph form
        reports its true covariance. The form $(I - KH)P^-$ would report $-1$.

  \item Take $P^- = \begin{bmatrix}10^4 & 9.99\\ 9.99 & 0.01\end{bmatrix}$ (correlation $0.999$), $H = [1\ \ 0]$ and $R = 10^{-4}$, an extremely accurate measurement of the
        first state, and compute everything in \code{float32}. The gain's first entry $10^4/(10^4 + 10^{-4})$ rounds to exactly $1$, so
        \[
          (I - KH)P^- = \begin{bmatrix}0 & 0\\ 9.5\times10^{-7} & 2.0\times10^{-5}\end{bmatrix},
        \]
        which is not symmetric, has a zero variance where the true value is $10^{-4}$, and whose symmetrized version has the eigenvalue
        $-1.1\times10^{-8} < 0$. The Joseph form gives $\begin{bmatrix}1.000\times10^{-4} & 9.99\times10^{-8}\\ 9.99\times10^{-8} & 1.999\times10^{-5}\end{bmatrix}$, positive definite and
        equal to the exact posterior to four digits. The first diagonal entry shows why: the Joseph form computes $(1 - K_1)^2P^-_{11} + K_1^2R = 0 + R$, so
        the measurement information enters as a positive term instead of as the tiny difference of two large ones.

  \item Only the covariances can be continued, since they do not depend on measurements. With $P_k^+ = 2P_k^-/(P_k^- + 2)$ and $P_{k+1}^- = P_k^+ + 1$:
        \begin{center}\small
        \begin{tabular}{@{}l c c c c c c c@{}}
        \toprule
        $k$ & 1 & 2 & 3 & 4 & 5 & 6 & 7 \\
        \midrule
        $P_k^-$ & 5 & 2.4286 & 2.0968 & 2.0236 & 2.0059 & 2.0015 & 2.0004 \\
        $K_k$ & 0.714 & 0.548 & 0.512 & 0.503 & 0.501 & 0.500 & 0.500 \\
        \bottomrule
        \end{tabular}
        \end{center}
        $P_k^- \to 2$, and the distance to the limit shrinks by a factor of about four per step: the derivative of the Riccati map at its fixed point is
        $\big(R/(P^-_\infty + R)\big)^2 = (2/4)^2 = \tfrac14$.

  \item At steady state $P = Pr/(P + r) + q$ for $P = P^-_\infty$. Multiplying by $P + r$: $P^2 + Pr = Pr + qP + qr$, i.e.\ $P^2 - qP - qr = 0$, whose positive root is
        \[
          P^-_\infty = \tfrac12\big(q + \sqrt{q^2 + 4qr}\big), \qquad K_\infty = \frac{P^-_\infty}{P^-_\infty + r}, \qquad P^+_\infty = \frac{P^-_\infty r}{P^-_\infty + r} = P^-_\infty - q .
        \]
        For $q/r \to 0$: $P^-_\infty \approx \sqrt{qr}$ and $K_\infty \approx \sqrt{q/r} \to 0$. The state barely moves, so the filter averages about $\sqrt{r/q}$ measurements and
        its error $\sqrt{qr}$ is much smaller than the sensor noise. For $q/r \to \infty$: $P^-_\infty \approx q + r$ and $K_\infty \to 1$. The state moves so much between
        measurements that each new one is nearly all the information; the filter follows the measurements, and $P^+_\infty \to r$.

  \item The implementation of \cref{alg:kf-update} on the model of \cref{fig:kf1d} (constant velocity, $\Dt = 0.1$~s, acceleration density $\sqrt{0.05}$,
        position noise variance 1, prior $\diag(4, 4)$) reproduces the figure. In the 200 steps of the figure's run no position error leaves the
        $\pm3\sigma$ envelope. One run says little, since the expected number is $200 \times 0.27\% \approx 0.5$; over 1000 independent runs the
        fraction is $0.29\%$, in agreement with the $0.27\%$ of a Gaussian with the correct variance.

  \item For the normalized innovations $\nu_k = r_k/\sqrt{S_k}$ of a long run (5000 steps, the first 50 discarded), the sample autocorrelation
        $\hat\rho(\ell) = \sum_k\nu_k\nu_{k+\ell}/\sum_k\nu_k^2$ should lie within $\pm1.96/\sqrt N \approx \pm0.028$ at every lag $\ell \ge 1$ if the innovations are white:
        \begin{center}\small
        \begin{tabular}{@{}l c c@{}}
        \toprule
        filter & $\hat\rho(1), \dots, \hat\rho(5)$ & mean NIS \\
        \midrule
        correct $Q$ & $0.00,\ 0.02,\ -0.01,\ 0.01,\ 0.02$ & 1.01 \\
        $Q/100$ & $0.44,\ 0.45,\ 0.43,\ 0.44,\ 0.44$ & 1.85 \\
        $100\,Q$ & $-0.11,\ -0.08,\ -0.09,\ -0.07,\ -0.05$ & 0.91 \\
        \bottomrule
        \end{tabular}
        \end{center}
        With too little process noise the filter trusts its model too much and lags behind the true motion, so consecutive innovations share the same
        sign (strong positive correlation, and an NIS well above 1). With too much process noise it overreacts to each measurement and partly
        corrects back at the next one (negative correlation). The whiteness test detects both, without ground truth.
\end{solutions}
